% ================================================================
%  sesion05.tex  —  Sesión 5 de 20
%  Rectas, segmentos y ángulos en la ciudad
%  Almería en Cifras y Formas · Matemáticas 1.º ESO
%  Compilar con:  pdflatex sesion05.tex  (dos pasadas)
% ================================================================
\documentclass[aspectratio=169, 10pt, spanish]{beamer}
\usepackage{almeria-beamer}

% ---------------------------------------------------------------
%  DATOS DE LA SESIÓN
% ---------------------------------------------------------------
\renewcommand{\SessionNumber}{05}
\renewcommand{\SessionTitle}{Rectas, segmentos y ángulos en la ciudad}
\renewcommand{\SessionName}{Sesión 05}
\renewcommand{\SessionObjective}{%
  Identificar y trazar puntos, rectas, segmentos y ángulos en planos
  urbanos, reconociendo rectas paralelas y perpendiculares en las
  calles de Almería.}
\renewcommand{\BlockName}{Bloque 2: Geometría urbana}

% ================================================================
\begin{document}

% ---------------------------------------------------------------
%  PORTADA
% ---------------------------------------------------------------
\begin{frame}[plain]
  \titlepage
\end{frame}

% ---------------------------------------------------------------
%  ÍNDICE DE LA SESIÓN
% ---------------------------------------------------------------
\begin{frame}{Índice de la sesión}
  \begin{columns}[t]
    \begin{column}{0.5\textwidth}
      \begin{enumerate}
        \item Punto, recta, semirrecta y segmento
        \item Ángulos y sus tipos
        \item Rectas paralelas y perpendiculares
        \item Almería: geometría en las calles
      \end{enumerate}
    \end{column}
    \begin{column}{0.5\textwidth}
      \begin{enumerate}
        \setcounter{enumi}{4}
        \item Diagrama completo de elementos
        \item Ejercicios multinivel (Bloom)
        \item Evidencia del día
        \item Error frecuente
        \item Resumen
      \end{enumerate}
    \end{column}
  \end{columns}
\end{frame}

% ---------------------------------------------------------------
\seccionframe{1. Punto, recta, semirrecta y segmento}
% ---------------------------------------------------------------

\begin{frame}{Elementos geométricos básicos (1/2)}
  \begin{columns}[T]
    \begin{column}{0.50\textwidth}
      \begin{definicionbox}{Punto}
        \AlmFontFootnote
        Posición exacta en el plano, \textbf{sin dimensiones}.\\
        Se representa con una letra mayúscula: $A$, $B$, $P$\ldots
      \end{definicionbox}
      \vspace{5pt}
      \begin{definicionbox}{Recta}
        \AlmFontFootnote
        Se extiende \textbf{infinitamente} en ambas direcciones;
        no tiene extremos. Se representa con una letra minúscula: $r$, $s$\ldots
      \end{definicionbox}
    \end{column}
    \begin{column}{0.46\textwidth}
      \begin{center}
        \begin{tikzpicture}[scale=0.88]
          % ---------- Punto ----------
          \node[punto] (P) at (0, 5.2) {};
          \node[font=\AlmFontSmall\bfseries, color=AlmAzulM, right=3pt of P] {$P$};
          \node[font=\AlmFontTiny, color=AlmGris, right=22pt of P] {punto};

          % ---------- Recta ----------
          \draw[recto, <->] (-1.8, 4.2) -- (2.2, 4.2);
          \node[font=\AlmFontTiny, color=AlmAzulM] at (-1.5, 4.45) {$r$};
          \node[font=\AlmFontTiny, color=AlmGris] at (2.8, 4.2) {recta};
        \end{tikzpicture}
      \end{center}
    \end{column}
  \end{columns}
\end{frame}

\begin{frame}{Elementos geométricos básicos (2/2)}
  \begin{columns}[T]
    \begin{column}{0.50\textwidth}
      \begin{definicionbox}{Semirrecta}
        \AlmFontFootnote
        Mitad de una recta: tiene \textbf{un extremo} (origen)
        y se extiende hacia un solo lado.
      \end{definicionbox}
      \vspace{5pt}
      \begin{definicionbox}{Segmento}
        \AlmFontFootnote
        Porción \textbf{finita} de recta entre dos puntos $A$ y $B$.
        Se escribe $\overline{AB}$. Tiene longitud medible.
      \end{definicionbox}
    \end{column}
    \begin{column}{0.46\textwidth}
      \begin{center}
        \begin{tikzpicture}[scale=0.88]
          % ---------- Semirrecta ----------
          \node[punto] at (-1.8, 3.2) {};
          \draw[recto, ->] (-1.8, 3.2) -- (2.2, 3.2);
          \node[font=\AlmFontTiny, color=AlmGris] at (2.8, 3.2) {semirrecta};

          % ---------- Segmento ----------
          \node[punto] (Sa) at (-1.8, 2.2) {};
          \node[punto] (Sb) at (2.2, 2.2) {};
          \draw[recto] (Sa) -- (Sb);
          \node[font=\AlmFontSmall\bfseries, color=AlmAzulM, below=2pt] at (Sa) {$A$};
          \node[font=\AlmFontSmall\bfseries, color=AlmAzulM, below=2pt] at (Sb) {$B$};
          \node[font=\AlmFontTiny, color=AlmGris] at (2.8, 2.2) {$\overline{AB}$};
          \draw[acota] (-1.8,1.78) -- (2.2,1.78);
          \node[font=\AlmFontTiny, color=AlmGris] at (0.2, 1.52) {longitud medible};

          \draw[AlmAmbar!60, line width=0.6pt, dashed]
            (-2.1,0) -- (3.2,0);
          \node[font=\AlmFontTiny\bfseries, color=AlmAmbar!80!black]
            at (0.5, -0.22) {elementos geométricos básicos};
        \end{tikzpicture}
      \end{center}
    \end{column}
  \end{columns}
\end{frame}

% ---------------------------------------------------------------
\seccionframe{2. Ángulos y sus tipos}
% ---------------------------------------------------------------

\begin{frame}{Qué es un ángulo}
  \begin{columns}[c]
    \begin{column}{0.48\textwidth}
      \begin{definicionbox}{Ángulo}
        \AlmFontSmall
        \textbf{Abertura} entre dos semirrectas
        que parten del mismo punto (vértice).\\
        Se mide en \textbf{grados} (°).
      \end{definicionbox}
      \vspace{8pt}
      \begin{notabox}[Cómo clasificarlo]
        \AlmFontFootnote
        Basta comparar su medida con 90° y 180°:
        antes de 90° es agudo, en 90° es recto,
        entre 90° y 180° es obtuso y en 180° es llano.
      \end{notabox}
    \end{column}
    \begin{column}{0.46\textwidth}
      {\AlmFontSmall
      \begin{tabular}{@{}lll@{}}
        \toprule
        \textbf{Tipo} & \textbf{Condición} & \textbf{Color}\\
        \midrule
        Agudo   & $0° < \alpha < 90°$
          & \textcolor{AlmAzulM}{azul}\\
        Recto   & $\alpha = 90°$
          & \textcolor{AlmAmbar!80!black}{ámbar}\\
        Obtuso  & $90° < \alpha < 180°$
          & \textcolor{AlmTerra}{rojo}\\
        Llano   & $\alpha = 180°$
          & \textcolor{AlmVerde}{verde}\\
        \bottomrule
      \end{tabular}}
      \vspace{8pt}
      \begin{ejemplobox}
        \AlmFontFootnote
        Una puerta poco abierta suele formar un
        \textbf{ángulo agudo}; la esquina de un folio
        es un \textbf{ángulo recto}.
      \end{ejemplobox}
    \end{column}
  \end{columns}
\end{frame}

\begin{frame}{Galería de tipos de ángulos}
  \begin{center}
    \begin{tikzpicture}[scale=0.86]
      % ---- Agudo ~55° ----
      \draw[recto, color=AlmAzulM]
        (0,0) -- (1.3,0);
      \draw[recto, color=AlmAzulM]
        (0,0) -- (0.74,1.06);
      \draw[AlmAzulM, line width=0.7pt]
        (0.45,0) arc (0:55:0.45);
      \node[font=\AlmFontTiny\bfseries, color=AlmAzulM] at (0.7,0.25)
        {$55°$};
      \node[font=\AlmFontTiny, color=AlmGris] at (0.65,-0.28)
        {agudo};

      % ---- Recto 90° ----
      \draw[recto, color=AlmAmbar!80!black]
        (2.5,0) -- (3.8,0);
      \draw[recto, color=AlmAmbar!80!black]
        (2.5,0) -- (2.5,1.3);
      \draw[rectoang]
        (2.5,0) ++(0.28,0) -- ++(0,0.28) -- ++(-0.28,0);
      \node[font=\AlmFontTiny\bfseries, color=AlmAmbar!80!black]
        at (3.2,0.45) {$90°$};
      \node[font=\AlmFontTiny, color=AlmGris] at (3.15,-0.28)
        {recto};

      % ---- Obtuso ~120° ----
      \draw[recto, color=AlmTerra]
        (4.8,0) -- (6.1,0);
      \draw[recto, color=AlmTerra]
        (4.8,0) -- (4.17,1.06);
      \draw[AlmTerra, line width=0.7pt]
        (5.22,0) arc (0:120:0.42);
      \node[font=\AlmFontTiny\bfseries, color=AlmTerra] at (5.05,0.35)
        {$120°$};
      \node[font=\AlmFontTiny, color=AlmGris] at (5.45,-0.28)
        {obtuso};

      % ---- Llano 180° ----
      \draw[recto, color=AlmVerde]
        (0,-1.5) -- (3,-1.5);
      \node[punto, fill=AlmVerde] at (1.5,-1.5) {};
      \draw[AlmVerde, line width=0.7pt]
        (1.92,-1.5) arc (0:180:0.42);
      \node[font=\AlmFontTiny\bfseries, color=AlmVerde] at (1.5,-1.1)
        {$180°$};
      \node[font=\AlmFontTiny, color=AlmGris] at (1.5,-1.82)
        {llano};
    \end{tikzpicture}
  \end{center}
  \vspace{6pt}
  \begin{center}
    \begin{notabox}[Idea clave]
      \AlmFontFootnote
      Cuanto mayor es la abertura, mayor es la medida del ángulo.
    \end{notabox}
  \end{center}
\end{frame}

% ---------------------------------------------------------------
\seccionframe{3. Rectas paralelas y perpendiculares}
% ---------------------------------------------------------------

\begin{frame}{Rectas paralelas y perpendiculares}
  \begin{columns}[T]
    \begin{column}{0.48\textwidth}
      \begin{definicionbox}{Rectas paralelas}
        \AlmFontSmall
        Dos rectas son \textbf{paralelas} si
        \textbf{nunca se cortan}, por mucho que
        las prolonguemos.\\[4pt]
        Siempre mantienen la \textbf{misma distancia}
        entre sí.\\[4pt]
        Notación: $r \parallel s$
      \end{definicionbox}
      \vspace{6pt}
      \begin{center}
        \begin{tikzpicture}[scale=0.75]
          \draw[recto, color=AlmAzulM]
            (0,0.5) -- (3.5,0.5)
            node[right, font=\AlmFontTiny\bfseries, color=AlmAzulM]{$r$};
          \draw[recto, color=AlmAzulM]
            (0,1.6) -- (3.5,1.6)
            node[right, font=\AlmFontTiny\bfseries, color=AlmAzulM]{$s$};
          % Marcas de paralelismo
          \draw[AlmAmbar, line width=1pt]
            (1.5,0.38) -- (1.5,0.62);
          \draw[AlmAmbar, line width=1pt]
            (1.5,1.48) -- (1.5,1.72);
          % Distancia
          \draw[acota] (2.8,0.5) -- (2.8,1.6);
          \node[font=\AlmFontTiny, color=AlmGris, right=2pt]
            at (2.8,1.05) {dist. cte.};
        \end{tikzpicture}
      \end{center}
    \end{column}
    \begin{column}{0.48\textwidth}
      \begin{definicionbox}{Rectas perpendiculares}
        \AlmFontSmall
        Dos rectas son \textbf{perpendiculares} si
        se cortan formando un \textbf{ángulo de 90°}
        (ángulo recto).\\[4pt]
        Notación: $r \perp s$
      \end{definicionbox}
      \vspace{6pt}
      \begin{center}
        \begin{tikzpicture}[scale=0.75]
          \draw[recto, color=AlmTerra]
            (0,1.05) -- (3.5,1.05)
            node[right, font=\AlmFontTiny\bfseries, color=AlmTerra]{$r$};
          \draw[recto, color=AlmTerra]
            (1.75,0) -- (1.75,2.3)
            node[above, font=\AlmFontTiny\bfseries, color=AlmTerra]{$s$};
          % Cuadradito ángulo recto
          \draw[rectoang]
            (1.75,1.05) ++(0.22,0) -- ++(0,0.22) -- ++(-0.22,0);
          \node[font=\AlmFontSmall\bfseries, color=AlmAmbar!80!black]
            at (2.35,1.5) {$90°$};
        \end{tikzpicture}
      \end{center}
    \end{column}
  \end{columns}
\end{frame}

% ---------------------------------------------------------------
\seccionframe{4. Almería: geometría en las calles}
% ---------------------------------------------------------------

\begin{frame}{Geometría en las calles de Almería}
  \begin{columns}[c]
    \begin{column}{0.50\textwidth}
      \begin{ejemplobox}
        \AlmFontSmall
        \textbf{Paralelas en Almería:}\\
        La \textbf{Avenida del Mediterráneo} y la
        \textbf{Calle del Obispo} discurren en la misma
        dirección; nunca se encuentran. Son \textbf{paralelas}.
        \vspace{4pt}\\
        \textbf{Perpendiculares en Almería:}\\
        La \textbf{Calle Reyes Católicos} corta a la
        \textbf{Rambla de Almería} formando un ángulo
        de aproximadamente \textbf{90°}. Son \textbf{perpendiculares}.
      \end{ejemplobox}
    \end{column}
    \begin{column}{0.46\textwidth}
      \begin{center}
        \begin{tikzpicture}[scale=0.85]
          % Fondo de manzana
          \fill[AlmFondo, rounded corners=3pt]
            (0,0) rectangle (4.4,3.8);
          % Calle horizontal 1 — Av. Mediterráneo (paralela)
          \fill[AlmAzulC!30]
            (0,2.6) rectangle (4.4,3.1);
          \node[font=\AlmFontTiny\bfseries, color=AlmAzul, align=center]
            at (2.2,2.85) {Av. Mediterráneo};
          % Calle horizontal 2 — C. del Obispo (paralela)
          \fill[AlmAzulC!30]
            (0,0.6) rectangle (4.4,1.1);
          \node[font=\AlmFontTiny\bfseries, color=AlmAzul, align=center]
            at (2.2,0.85) {C. del Obispo};
          % Marcas de paralelismo
          \draw[AlmAmbar, line width=1.2pt]
            (2.0,2.62) -- (2.0,3.08);
          \draw[AlmAmbar, line width=1.2pt]
            (2.0,0.62) -- (2.0,1.08);
          % Calle vertical — C. Reyes Católicos (perpendicular a Rambla)
          \fill[AlmTerra!25]
            (1.8,0) rectangle (2.3,3.8);
          \node[font=\AlmFontTiny\bfseries, color=AlmTerra,
                rotate=90, align=center]
            at (2.05,1.85) {C. Reyes Católicos};
          % Cuadradito ángulo recto en intersección Rambla
          \draw[rectoang, draw=AlmTerra, line width=0.9pt]
            (2.3,1.1) ++(0,0) -- ++(0.2,0) -- ++(0,0.2);
          % Leyenda
          \node[font=\AlmFontTiny, color=AlmAzulM] at (3.6,3.45)
            {\textcolor{AlmAzulM}{$\parallel$} paralelas};
          \node[font=\AlmFontTiny, color=AlmTerra] at (3.65,3.15)
            {\textcolor{AlmTerra}{$\perp$} perpendicular};
        \end{tikzpicture}
      \end{center}
    \end{column}
  \end{columns}
\end{frame}

\begin{frame}{Las calles son segmentos}
  \begin{columns}[c]
    \begin{column}{0.56\textwidth}
      \begin{notabox}[Idea clave]
        \AlmFontSmall
        Cualquier calle tiene un inicio y un final:
        en el plano urbano la representamos como un
        \textbf{segmento}, no como una recta infinita.
      \end{notabox}
      \vspace{8pt}
      \begin{ejemplobox}
        \AlmFontFootnote
        Si una calle va desde la plaza A hasta el cruce B,
        podemos escribirla como $\overline{AB}$ y medir
        su longitud aproximada.
      \end{ejemplobox}
    \end{column}
    \begin{column}{0.38\textwidth}
      \begin{center}
        \begin{tikzpicture}[scale=0.9]
          \node[punto] (A) at (0,0) {};
          \node[punto] (B) at (3.2,0) {};
          \draw[recto] (A) -- (B);
          \node[font=\AlmFontSmall\bfseries, color=AlmAzulM, below=2pt] at (A) {$A$};
          \node[font=\AlmFontSmall\bfseries, color=AlmAzulM, below=2pt] at (B) {$B$};
          \draw[acota] (0,-0.45) -- (3.2,-0.45);
          \node[font=\AlmFontTiny, color=AlmGris] at (1.6,-0.72) {tramo de calle};
        \end{tikzpicture}
      \end{center}
    \end{column}
  \end{columns}
\end{frame}

% ---------------------------------------------------------------
\seccionframe{5. Diagrama completo de elementos}
% ---------------------------------------------------------------

\begin{frame}{Todos los elementos en un mapa urbano}
  \framesubtitle{Plano esquemático de un barrio de Almería}
  \begin{center}
    \begin{tikzpicture}[scale=1.0]
      % Fondo
      \fill[AlmFondo, rounded corners=4pt]
        (-0.6,-0.6) rectangle (9.6,5.0);

      % ---- Calles paralelas horizontales ----
      \draw[recto, color=AlmAzulM, line width=1.8pt]
        (0,4.0) -- (9,4.0);
      \draw[recto, color=AlmAzulM, line width=1.8pt]
        (0,2.2) -- (9,2.2);
      \draw[recto, color=AlmAzulM, line width=1.8pt]
        (0,0.4) -- (9,0.4);

      % Marcas de paralelismo en las tres rectas
      \foreach \y in {4.0,2.2,0.4}{
        \draw[AlmAmbar, line width=1.0pt]
          (4.3,\y-0.12) -- (4.3,\y+0.12);
        \draw[AlmAmbar, line width=1.0pt]
          (4.6,\y-0.12) -- (4.6,\y+0.12);
      }
      \node[font=\AlmFontTiny\bfseries, color=AlmAzulM] at (-0.3,4.0)
        {$r_1$};
      \node[font=\AlmFontTiny\bfseries, color=AlmAzulM] at (-0.3,2.2)
        {$r_2$};
      \node[font=\AlmFontTiny\bfseries, color=AlmAzulM] at (-0.3,0.4)
        {$r_3$};
      \node[font=\AlmFontTiny, color=AlmAzulM] at (9.3,4.0) {Av.\ Mar};
      \node[font=\AlmFontTiny, color=AlmAzulM] at (9.3,2.2) {Rambla};
      \node[font=\AlmFontTiny, color=AlmAzulM] at (9.3,0.4) {C.\ Obispo};

      % ---- Calle perpendicular ----
      \draw[rectot, line width=1.8pt]
        (2.8,-0.4) -- (2.8,4.8);
      \node[font=\AlmFontTiny\bfseries, color=AlmTerra, rotate=90]
        at (2.55,2.2) {C.\ Reyes Católicos};
      % Cuadradito ángulo recto
      \draw[rectoang, draw=AlmTerra, line width=0.9pt]
        (2.8,2.2) ++(0.18,0) -- ++(0,0.18) -- ++(-0.18,0);

      % ---- Segmentos verdes ----
      \draw[rectov, line width=2pt]
        (5.5,0.4) -- (5.5,2.2);
      \node[font=\AlmFontTiny\bfseries, color=AlmVerde, rotate=90]
        at (5.3,1.3) {$\overline{CD}$};
      \draw[rectov, line width=2pt]
        (7.0,2.2) -- (8.5,4.0);
      \node[font=\AlmFontTiny\bfseries, color=AlmVerde]
        at (8.0,3.35) {$\overline{EF}$};
      \draw[rectov, line width=2pt]
        (1.0,0.4) -- (2.8,0.4);
      \node[font=\AlmFontTiny\bfseries, color=AlmVerde]
        at (1.9,0.62) {$\overline{AB}$};

      % ---- Puntos notables ----
      \node[punto] (Pt) at (2.8,4.0) {};
      \node[font=\AlmFontTiny\bfseries, color=AlmAzulM, above right=0pt]
        at (Pt) {$P$};
      \node[punto] (Qt) at (5.5,2.2) {};
      \node[font=\AlmFontTiny\bfseries, color=AlmAzulM, above right=0pt]
        at (Qt) {$Q$};

      % ---- Ángulo agudo entre segmento oblicuo y horizontal ----
      \draw[AlmAmbar!80!black, line width=0.7pt]
        (7.28,2.2) arc (0:50:0.28);
      \node[font=\AlmFontTiny\bfseries, color=AlmAmbar!80!black]
        at (7.7,2.45) {agudo};

      % ---- Leyenda ----
      \node[draw=AlmAzulC!50, fill=AlmBlanco, rounded corners=2pt,
            inner sep=4pt, font=\AlmFontTiny, align=left]
        at (1.5,3.35) {
          \textcolor{AlmAzulM}{—} paralelas ($r_1 \parallel r_2 \parallel r_3$)\\
          \textcolor{AlmTerra}{—} perpendicular ($\perp$)\\
          \textcolor{AlmVerde}{—} segmentos ($\overline{AB}$, $\overline{CD}$, $\overline{EF}$)
        };
    \end{tikzpicture}
  \end{center}
\end{frame}

% ---------------------------------------------------------------
\seccionframe{6. Ejercicios multinivel — Bloom}
% ---------------------------------------------------------------

\begin{frame}{Ejercicio RECORDAR}
  \bloomtag{RECORDAR}\quad\dificultad{1}
  \vspace{4pt}
  \begin{recordarbox}
    \AlmFontFootnote
    Con \textbf{regla y compás} (o solo regla), dibuja en
    tu cuaderno cada uno de los siguientes elementos:
    \begin{enumerate}
      \item Una \textbf{recta} $r$ (indica la doble flecha
            para mostrar que no tiene extremos).
      \item Un \textbf{segmento} $\overline{AB}$ de 5 cm.
            Mide y anota la longitud.
      \item Un \textbf{ángulo recto} de 90°.
            Marca el cuadradito del ángulo recto.
      \item Un \textbf{ángulo agudo} de aproximadamente 60°.
      \item Un \textbf{ángulo obtuso} de aproximadamente 110°.
    \end{enumerate}
    \vspace{4pt}
    Escribe junto a cada dibujo el nombre del elemento
    y su definición en una frase breve.
  \end{recordarbox}
\end{frame}

\begin{frame}{Ejercicio COMPRENDER}
  \bloomtag{COMPRENDER}\quad\dificultad{2}
  \vspace{4pt}
  \begin{comprenderbox}
    \AlmFontFootnote
    Observa una fotografía del centro histórico de Almería
    (o el plano que facilita el profesor).\\[6pt]
    Identifica y \textbf{justifica} con una frase cada uno:
    \begin{enumerate}
      \item \textbf{Dos calles paralelas.}\\
            Justificación: «Son paralelas porque\ldots»
      \item \textbf{Una intersección perpendicular.}\\
            Justificación: «Son perpendiculares porque\ldots»
      \item \textbf{Un ángulo obtuso} visible en un cruce o edificio.\\
            Justificación: «Es obtuso porque mide más de 90°\ldots»
    \end{enumerate}
    \vspace{4pt}
    \begin{notabox}[Pista]
      \AlmFontFootnote
      Para justificar el paralelismo: comprueba que
      las calles \emph{nunca} se cruzan en el mapa.
    \end{notabox}
  \end{comprenderbox}
\end{frame}

\begin{frame}{Ejercicio APLICAR}
  \bloomtag{APLICAR}\quad\dificultad{3}
  \vspace{4pt}
  \begin{aplicarbox}
    \AlmFontFootnote
    En el \textbf{plano simplificado de Almería} que te
    entrega el profesor:
    \begin{itemize}
      \item Traza con \textcolor{AlmAzulM}{\textbf{azul}}
            las \textbf{2 calles paralelas} que identifiques.
            Añade las marcas de paralelismo
            (\,$\shortmid$\;$\shortmid$\,) sobre ellas.
      \item Traza con \textcolor{AlmTerra}{\textbf{rojo}}
            la \textbf{intersección perpendicular} que
            identifiques. Dibuja el cuadradito de 90°.
      \item Marca con \textcolor{AlmVerde}{\textbf{verde}}
            \textbf{3 segmentos} que representen tramos de calles.
            Escribe su nombre ($\overline{AB}$, $\overline{CD}$,
            $\overline{EF}$) y mide su longitud aproximada en cm.
      \item Señala \textbf{2 ángulos agudos} y \textbf{1 ángulo obtuso}
            que aparezcan en algún cruce o edificio del plano.
    \end{itemize}
  \end{aplicarbox}
\end{frame}

% ---------------------------------------------------------------
%  EVIDENCIA
% ---------------------------------------------------------------
\begin{frame}{Evidencia del día}
  \begin{evidenciabox}
    \AlmFontFootnote
    Entrega tu \textbf{Plano anotado de Almería} con:
    \begin{itemize}
      \item \textcolor{AlmAzulM}{\textbf{Paralelas en azul}}
            con marcas de paralelismo y nombre de las calles.
      \item \textcolor{AlmTerra}{\textbf{Perpendiculares en rojo}}
            con el cuadradito de ángulo recto dibujado.
      \item \textcolor{AlmVerde}{\textbf{Segmentos en verde}}
            con sus extremos nombrados y longitud aproximada.
      \item Al menos \textbf{2 ángulos} etiquetados con su tipo
            (agudo / recto / obtuso / llano).
    \end{itemize}
  \end{evidenciabox}
\end{frame}

\begin{frame}{Evidencia del día: tabla resumen}
  \begin{evidenciabox}
    \AlmFontFootnote
    Añade al final una \textbf{tabla resumen} como esta:
    \vspace{6pt}
    \begin{center}
      \begin{tabular}{@{}lll@{}}
        \toprule
        \textbf{Elemento} & \textbf{Ejemplo en Almería} & \textbf{Justificación}\\
        \midrule
        Paralelas & Av.\,Mediterráneo y C.\,Obispo & Nunca se cortan\\
        Perpendiculares & C.\,Reyes Católicos $\perp$ Rambla & Forman 90°\\
        Segmento & Tramo $\overline{AB}$ & Tiene inicio y final\\
        \bottomrule
      \end{tabular}
    \end{center}
  \end{evidenciabox}
\end{frame}

% ---------------------------------------------------------------
%  ERROR FRECUENTE
% ---------------------------------------------------------------
\begin{frame}{¡Cuidado con este error!}
  \begin{errorbox}
    \begin{columns}[c]
      \begin{column}{0.50\textwidth}
        \incorrecto\ \textbf{Error frecuente}\\[6pt]
        Llamar «recta» a una calle porque parece
        infinitamente larga.\\[6pt]
        \begin{center}
          \cancel{«La Calle Mayor es una \textbf{recta}»}
        \end{center}
        \vspace{4pt}
        \AlmFontFootnote
        Toda calle tiene un inicio y un final:
        \textbf{siempre es un segmento}.
        Una recta matemática no tiene extremos
        y se extiende al infinito en ambos sentidos,
        lo que es imposible en el mundo real.
      \end{column}
      \begin{column}{0.46\textwidth}
        \correcto\ \textbf{Correcto}\\[6pt]
        \AlmFontSmall
        «La Calle Mayor es el
        \textbf{segmento} $\overline{AB}$,
        donde $A$ es su inicio y $B$ su final.»\\[8pt]
        \begin{center}
          \begin{tikzpicture}[scale=0.8]
            \node[punto] (A) at (0,0) {};
            \node[font=\AlmFontSmall\bfseries, color=AlmAzulM,
                  below=2pt] at (A) {$A$};
            \node[punto] (B) at (3.5,0) {};
            \node[font=\AlmFontSmall\bfseries, color=AlmAzulM,
                  below=2pt] at (B) {$B$};
            \draw[recto] (A) -- (B);
            \draw[acota] (0,-0.42) -- (3.5,-0.42);
            \node[font=\AlmFontTiny, color=AlmGris] at (1.75,-0.65)
              {longitud finita};
          \end{tikzpicture}
        \end{center}
      \end{column}
    \end{columns}
  \end{errorbox}
\end{frame}

\begin{frame}{Regla de oro}
  \begin{notabox}[Para no confundirlo]
    \AlmFontSmall
    \textbf{Recta} $\to$ sin extremos, infinita.\\[4pt]
    \textbf{Segmento} $\to$ dos extremos, longitud medible.
  \end{notabox}
  \vspace{10pt}
  \begin{center}
    \begin{tikzpicture}[scale=0.9]
      \draw[recto,<->] (-1.8,0.8) -- (1.8,0.8);
      \node[font=\AlmFontTiny, color=AlmGris] at (0,1.1) {recta};
      \node[punto] (A) at (-1.4,-0.6) {};
      \node[punto] (B) at (1.4,-0.6) {};
      \draw[recto] (A) -- (B);
      \node[font=\AlmFontTiny, color=AlmGris] at (0,-1.0) {segmento};
    \end{tikzpicture}
  \end{center}
\end{frame}

% ---------------------------------------------------------------
%  RESUMEN
% ---------------------------------------------------------------
\begin{frame}{Resumen de la sesión 05}
  \begin{resumenbox}
    \begin{itemize}
      \item \textbf{Punto}: posición sin dimensiones. \textbf{Recta}: sin extremos, infinita. \textbf{Semirrecta}: un extremo. \textbf{Segmento}: dos extremos, longitud medible.
      \item Un \textbf{ángulo} es la abertura entre dos semirrectas
            con el mismo origen; se mide en grados (°).
      \item Tipos: \textbf{agudo} ($< 90°$),
            \textbf{recto} ($= 90°$),
            \textbf{obtuso} ($90°$--$180°$),
            \textbf{llano} ($= 180°$).
      \item \textbf{Paralelas}: nunca se cortan ($r \parallel s$).
      \item \textbf{Perpendiculares}: se cortan en 90° ($r \perp s$).
      \item Las calles de Almería son \textbf{segmentos},
            no rectas, porque tienen inicio y final.
    \end{itemize}
  \end{resumenbox}
  \vspace{6pt}
  \begin{center}
    {\AlmFontSmall\color{AlmAzulM}
    \faArrowCircleRight\enskip
    \textbf{Próxima sesión:}
    Polígonos en la ciudad — triángulos y cuadriláteros en Almería.}
  \end{center}
\end{frame}

\end{document}



